\chapter{Introduction}\label{ch:introduction}

% --- §1.1 and §1.2 written 2026-09-28.  Drafted in the chat as prose over three passes and
%   approved before this write (CLAUDE.md §2, steps 5-6).  Snapshot:
%   chapter_01_introduction.pre1112.bak.tex.  Budget: §1.1 2.5 pp max, §1.2 1 p max (Abram);
%   ~965 and ~365 words, page counts from his build.
%
% WHAT THESE TWO SECTIONS ARE FOR.  §15 O-20 (b): the ten terms §1.3 uses must be established
%   before §1.3 uses them -- Göbel's own condition for naming answer retrieval in a research
%   question.  All ten land here.  §1.1: retrieved context, generator, structure-aware formula
%   index, formula search, answer retrieval (contrasted against formula search, as she asked),
%   formula representation, text-only retrieval.  §1.2: multi-index retriever, retrieval quality,
%   faithfulness, factual correctness, generation without retrieval.
%
% INTRODUCED HERE, DEFINED IN CHAPTER 3 -- the boundary this write turns on.  Chapter 1 gives each
%   term a plain gloss and Chapter 3 keeps every definition site: answer retrieval and formula
%   search are §3.1's, the tree representations §3.2's, faithfulness and factual correctness
%   §3.7's.  Italics follow that split rather than the usual first-use rule (CLAUDE.md §3), because
%   italicising a term in both places would give it two definition sites, which is the defect the
%   first-draft review named.
%   ⚠ FOUR TERMS ARE ITALICISED ANYWAY, on Abram's instruction 2026-09-28: *multi-index retriever*,
%   *faithfulness*, *factual correctness*, *generation without retrieval*.  The first and the last
%   are this thesis's own names and have no other site, so Chapter 1 is where they belong.  The
%   middle two ARE italicised again in §3.7, deliberately -- a one-clause gloss here, the
%   construction there.  §7a should carry a row for each recording the split.  The five plain terms
%   in §1.1 could take the same emphasis if the asymmetry reads badly on the page; raised and left.
%
% NO COINED TERM AND NO FIGURE.  Nothing from §5.2.2 or §6.2 -- not blocked, flat, formula block,
%   fanout, floor, configuration or judge pool -- and no number anywhere, so neither section can go
%   stale against the registry.  The sparse and dense indexes are described, never named; so are
%   the trees.  "formulas", not "formulae".
%
% THE THREE DIFFICULTIES ARE NUMBERED IN THE PROSE (Abram, 2026-09-28): the arrangement on the
%   retrieval side, whether a better ranking reaches the answer, and whether any of it can be
%   measured on a collection judged beforehand.  That is §2.5's three-gap shape, so §1.1 names the
%   difficulties and §2.5 names what the literature has and has not done about each.  Without the
%   ordinals the paragraphs read as a list of unrelated worries -- his note on the draft.
%
% ⚠ §1.1's LAST-BUT-ONE PARAGRAPH IS §15 O-6's ANSWER: the title says *scientific text* and the
%   evaluation is community mathematical question answering.  It concedes the gap and scopes every
%   claim to the collection.  THE TITLE ITSELF IS STILL UNDECIDED -- O-6's other half.
%
% CITATIONS, all already in sources.bib.  lewis2020rag and gao2024ragsurvey for the arrangement and
%   what it usually searches; dadure2024mirreview for the field having known the problem from the
%   start; davila2017tangents and mansouri2019tangentcft for indexing structure; mansouri2022arqmath3
%   for formula search being where such indexes are compared; zanibbi2020arqmath for the collection
%   and for answer retrieval; salemi2024erag, henkel2024mathrag and alshammari2026mathnet for the
%   published findings disagreeing; buckley2004incomplete for pooled judgments.
%   ⚠ alshammari2026mathnet is one of §15 O-22's seven entries with TODO fields and no PDF.  §2.3
%   already leans on it, so this adds no new debt, but it is the one claim in these two sections
%   resting on an abstract.  The three-key cluster is deliberate: the claim IS that they disagree.
%
% Cut on Abram's reading of the drafts, each for a reason worth keeping:
%   - "repair" twice for what RAG and structure indexing do -> "the standard approach" and "the
%     solution"; plainer, and "repair" read as though the field had broken something.
%   - "Only one of them reaches an ordinary text index" -- the next sentence says it better.
%   - "Keep the symbols and rearrange them and the index sees much the same formula, though it now
%     computes something else" -- the rename example already carries the point.
%   - "built alike rather than written alike" -> "mean the same thing rather than ... written the
%     same way", which is the property the index is after.
%   - "an index apiece" -> "one index per representation".
%   - "One boundary belongs at the outset" -> "The scope of what is claimed here should be clear
%     from the start".
%   - the optional scope paragraph in §1.2, carrying the three §1c-i scope decisions.  ⚠ THEY ARE
%     THEREFORE STATED NOWHERE IN THE THESIS: §1.6 cannot hold them at 173 words and §7.2 was
%     written without them.  Manifest §4 and §10 still expect them.  Decide before submission.
%
% "Mathematics Stack Exchange", the site's own name, in all three places it appears in the thesis --
%   §1.1, §2.1 and §3.1.  §3.1 said "Math Stack Exchange" until 2026-09-28, when Abram had it swept;
%   its paragraph was reflowed with the change.  Bib titles keep whatever their papers print.
\section{Motivation}\label{sec:intro-motivation}

A student with a mathematical question asks it in words and formulas together. A language model can
answer such a question from what it learnt in training, and it will answer whether or not it knows.
A standard approach is to put a search step in front of it, so that the model is provided with the
relevant context. The model then writes its answer from the question and that
context~\cite{lewis2020rag}. Text is the most widely used retrieval source in
systems built this way~\cite{gao2024ragsurvey}.

Mathematics is not prose. A post on a mathematics site carries its argument in two registers at
once, the words that explain and the formulas that do the work, and a retriever has to read both. A
formula enters an ordinary text index as the characters it was typed with, so it can be matched only
where those same characters appear. Rename the variables and the index sees an unrelated string.
The index never holds the arrangement itself --- operators and their operands, and the spatial
arrangement of variables and operators. This is where the meaning of a formula lies. This is the
problem mathematical information retrieval has worked on since it began~\cite{dadure2024mirreview}.

The solution is to index the structure rather than the characters. A formula can be parsed into a
tree, and that tree turned into something an index can hold and compare, so that two formulas count
as similar when they mean the same thing rather than when they are written the same
way~\cite{davila2017tangents,mansouri2019tangentcft}. An index of that kind is a structure-aware
formula index. These indexes work when the query is itself a formula. Formula search --- given one
formula, return the formulas most like it --- is what they were designed for, and where they are
compared against one another~\cite{mansouri2022arqmath3}.

The task this thesis studies is the other one. Answer retrieval takes a mathematical question,
stated in words and formulas, and returns the posts that answer it~\cite{zanibbi2020arqmath}. What
separates it from formula search is not its input but its standard: an answer is relevant when it
resolves the question, so a post whose formulas match the question symbol for symbol can still be
worthless, and a restatement of the question is the clearest case of it. A formula index that finds
similar formulas well has therefore not been shown to find answers well. The two cannot be read off
one another.

The first difficulty follows from that, and it is one of arrangement. A formula can be represented
as a tree in more than one way, and each of these formula representations keeps something the others
drop, so a system can hold several of them at once, one index per representation, rather than
choosing between them. Nor can a formula index work alone: the words around a formula are what
separate an answer deriving a result from a comment repeating the question, so its ranking has to be
combined with that of text-only retrieval --- retrieval over the words of a post and nothing else.
How several formula representations behave as a group, and how that group meets text retrieval, is
underexplored where answers are what is retrieved, as Chapter~\ref{ch:related-work} sets out.

The second difficulty is that a better ranking is not the end of the matter. What a reader sees is
the answer the generator wrote, and a ranking reaches that answer only through the context it
supplied. Those are two measurements, and whether the first carries into the second is unsettled:
the published findings
disagree~\cite{salemi2024erag,henkel2024mathrag,alshammari2026mathnet}, and none of that work
varies how formulas are represented in the retriever.

The third difficulty belongs to the measuring rather than to the system. A test collection is judged
once, on what the systems taking part at the time retrieved; everything outside that is left
unjudged and normally counted as not relevant~\cite{buckley2004incomplete}. A system built
afterwards retrieves posts nobody ever looked at, and is marked down for finding them. Nothing can
be claimed about such a system until the instruments that score it have been shown to work, which is
why two of them are established here before any result is reported.

The scope of what is claimed here should be clear from the start. The problem described above is
general: wherever a document carries mathematics in its text, an index reading only the characters
loses the structure, and the context a generator receives is the poorer for it. What this thesis
measures is narrower. Every experiment runs on one collection of community mathematical question
answering, built from a snapshot of Mathematics Stack Exchange~\cite{zanibbi2020arqmath}. The
questions are asked by people who want help with a problem rather than written up for publication,
the answers are posts rather than articles, and no scientific paper is retrieved anywhere in this
thesis. Every claim made here is a claim about that collection.

What remains is a system to build and three things to find out about it.
Section~\ref{sec:intro-objectives} states them.

\section{Objectives}\label{sec:intro-objectives}

The objective is to build a retriever for answer retrieval that reads formula structure, and to find
out what it changes --- first in the ranking, and then in the answer a generator writes from it.

Three structure-aware formula indexes are built, one per formula representation, so that the
representations are compared rather than assumed. They are placed beside two text indexes, one
matching the words of a post literally and one matching them by meaning. A question is put to all of
them, and the rankings they return are merged into one. A retriever that combines a text index with
at least one formula index is a \emph{multi-index retriever}. How the merge is done is not given in
advance, so it is varied: which rankings enter it, and whether the three formula rankings enter it
together or separately.

Each such retriever is scored against relevance judgments for the same questions, and that score is
its retrieval quality. Whether adding formula indexes to text-only retrieval raises it is the first
thing to settle.

The second is what that reaches. The best of the multi-index retrievers hands the top of its ranking
to one generator as the retrieved context, and the answers are scored twice: for how much of what an
answer says its context supplied, which is \emph{faithfulness}, and for how much of what it says is
correct, which is \emph{factual correctness}. Both readings are set against answers the same
generator writes for the same questions with no context at all --- \emph{generation without
retrieval} --- and against the answers built from text-only retrieval's context. The third thing to
settle is whether any of it holds when the generator is larger.

Two things are outside the objective. The retriever is assembled from published components and none
of them is trained further. One prompt is written and held fixed, so that what separates two sets of
answers is the context and nothing else.

Section~\ref{sec:intro-rqs} puts all of this as questions.

% --- §1.3 and §1.4 rewritten 2026-09-24 against manifest §1b and §1c, whose RQ and hypothesis
% sets were restated that day. Drafted in chat as prose and approved by Abram before this write
% (CLAUDE.md §2, steps 5-6). Snapshot: chapter_01_introduction.preRQ0924.bak.tex.
%
% Abram's decisions, 2026-09-24, each a deliberate departure from the old file:
%   - DECISION RULES are NOT stated per hypothesis here. They live in manifest §1c and are applied
%     at the verdict sites. Manifest §5 row 1.4 updated to match.
%   - THE TWO SELECTION RULES ("the strongest text-only retrieval", "the best multi-index
%     retriever" = highest ndcg'@10 under the judge's labels within each group, 78 topics) are NOT
%     stated here; manifest §1b keeps them.
%   - The ANSWERING-EFFICIENCY concession is dropped from Chapter 1; §7.2 carries it (manifest §1b).
%
% Retired here: old H1.1 (transfer) and old H1.2 (containment) carry no verdict since 2026-09-24
% and are narrative in Chapter 5. That also removes Chapter 1's only use of *formula block*
% (closed §15 O-14 d). Renumbering: old H1.3 -> H1.1 and H1.2(a); old H1.4 -> H1.3;
% old H1.5 -> H1.2(b) and H1.4; old H2.2/H2.3 -> H2.3/H2.4; old H2.4 -> H3.1. H2.2 is new.
% The thesis never mentions any renumbering -- its reader never saw the old IDs.
%
% ✅ DONE 2026-09-28: §1.1 and §1.2 establish all ten terms this section uses -- answer retrieval,
%   structure-aware formula index, text-only retrieval, multi-index retriever, factual correctness,
%   faithfulness, retrieval quality, retrieved context, generation without retrieval, generator --
%   with *answer retrieval* set against formula search, Göbel's own condition for naming it in a
%   research question (§15 O-20 b, closed). See the note above §1.1 for which term is introduced
%   where, and for the introduced-here / defined-in-Chapter-3 split.
% ✅ DONE 2026-09-28 (Abram): every hypothesis now points at the section that tests it, so the list
%   is parallel across the two groups and the opening paragraph names the two ruling sites.
%   Retrieval, repointed after the §15 O-21 split: H1.1 -> §5.4.3 `sec:res-add-one` (was
%   `sec:res-single-index`, §5.4.2, which decides no hypothesis), H1.2 -> §5.4.4
%   `sec:res-add-three`, H1.3 -> §5.4.5 `sec:res-fusion` (already correct), H1.4 -> §5.4.6
%   `sec:res-ablation`.  The keys are semantic, so they followed the move and the numbers are
%   LaTeX's.
%   Generation: H2.1 and H2.2 -> §6.3.2 `sec:gen-res-context`, H2.3 and H2.4 -> §6.3.3
%   `sec:gen-res-tau`, H3.1 -> §6.3.4 `sec:gen-res-sizes`, replacing five bare
%   Chapter~\ref{ch:generation} pointers.  ⚠ The old plan in manifest §9b sent H2.1 to §6.3.1 and
%   H2.2 to §6.3.2; that mapping predates the §6.3 OVERVIEW subsection, which took 6.3.1, and
%   §6.3.2 opens by naming both hypotheses, so both point there.  Manifest §9b corrected with this
%   write.  Two hypotheses sharing a pointer is the existing pattern -- H2.3 and H2.4 were always
%   planned to share one.
% NOTE If Göbel holds to her earlier ruling that generator size carries no hypothesis
%   (§15 O-2; checklist G4), RQ3 and H3.1 come out and the ladder is reported with no verdict.
%
% --- §1.4 OPENING PARAGRAPH EDITED 2026-09-28, on Abram's instruction, two changes.
%   (i) "the generation hypotheses at the end of Chapter 6" -> Section~\ref{sec:gen-verdicts}.
%       §6.4 exists and is labelled since 2026-09-26, so the vague pointer is no longer needed and
%       both groups' ruling sites are now named the same way.
%   (ii) THE 0.02 SENTENCE IS DELETED -- "Throughout, two scores differing by 0.02 or less are
%       reported as indistinguishable rather than as one exceeding the other." His instruction, and
%       it was the one sentence in the thesis that gave the limit CHAPTER-WIDE scope: the limit is
%       the relevance judge's and governs retrieval scores alone (§4.4, §5.1), and §6.1 says in
%       terms that generation scores are read as reported with no equivalent limit. "Throughout"
%       therefore claimed more than §4.4 measured. ⚠ Chapter 1 now states NO reading rule at all;
%       §5.1 states it for retrieval and §4.4 in full, so no verdict loses its basis, and §15
%       O-19 (b)'s "Chapter 1 needs no edit" is answered by deletion rather than by rewording.
%   The five generation hypotheses still carry \hfill Chapter~\ref{ch:generation} rather than a
%   §6.3.x pointer -- the TODO above stands, and it is a separate write.

\section{Research questions}\label{sec:intro-rqs}

Three questions organise the work. The first asks whether structure-aware formula indexes improve
answer retrieval over text alone, and whether the way they are combined governs the result. The
second asks how much the retrieved context of the resulting retriever improves the answers a
generator produces, and how faithful those answers are to it. The third checks that the second's
answers do not depend on the size of the generator. Retrieval and generation are kept apart
deliberately, each answered in a chapter of its own, so that retrieval is settled before generation
is opened.

\begin{description}[leftmargin=3.2em,style=nextline]
  \item[\textbf{RQ1}] Does adding structure-aware formula indexes to text-only retrieval improve
    answer retrieval on mathematical questions, and does the way those indexes are combined govern
    the effectiveness of the resulting multi-index retriever?
    \hfill Chapter~\ref{ch:retrieval}
\end{description}

\begin{description}[leftmargin=5.6em,style=nextline,topsep=0.2em]
  \item[\textbf{RQ1.1}] Does adding one formula index to the strongest text-only retrieval improve
    answer retrieval?
  \item[\textbf{RQ1.2}] Does adding all three formula indexes to the strongest text-only retrieval
    improve answer retrieval, and improve on adding one, and does the way they are combined make a
    difference?
\end{description}

\begin{description}[leftmargin=3.2em,style=nextline]
  \item[\textbf{RQ2}] To what extent does the retrieved context of the best multi-index retriever
    of RQ1 improve the factual correctness of generated answers over generation without retrieval
    and over the strongest text-only retrieval, how faithful are those answers to that context, and
    do both rise with the quality of the retrieval that supplied it?
    \hfill Chapter~\ref{ch:generation}
\end{description}

\begin{description}[leftmargin=5.6em,style=nextline,topsep=0.2em]
  \item[\textbf{RQ2.1}] Does the context of the best multi-index retriever of RQ1 make answers more
    factually correct than generation without retrieval, and more factually correct and more
    faithful than the strongest text-only retrieval?
  \item[\textbf{RQ2.2}] Across all the retrieval systems compared in RQ1, do factual correctness and
    faithfulness rise with retrieval quality?
\end{description}

\begin{description}[leftmargin=3.2em,style=nextline]
  \item[\textbf{RQ3}] Do the answers to RQ2 hold as the generator grows from 1.5 to 7~billion
    parameters?
    \hfill Chapter~\ref{ch:generation}
\end{description}

Two further conditions have to hold before any answer to these questions can be trusted, and
neither is a research question in its own right. Both are established in the experimental setup.

\begin{description}[leftmargin=3.2em,style=nextline]
  \item[\textbf{P1}] The evaluation pipeline reproduces the organisers' published figures on their
    own baseline run (Section~\ref{sec:met-calibration}).
  \item[\textbf{P2}] Retrieval built after the test collection was assembled can nonetheless be
    scored on a judgment pool it never contributed to --- through hole-rate reporting, prime
    metrics and a relevance judge validated before it is relied upon
    (Section~\ref{sec:met-judge-validation}).
\end{description}

\section{Hypotheses}\label{sec:intro-hypotheses}

The hypotheses divide as the questions do, into one group for retrieval and one for generation.
Each is stated as a positive prediction that the data is free to reject. The retrieval hypotheses
are ruled in Section~\ref{sec:ret-verdicts} and the generation hypotheses in
Section~\ref{sec:gen-verdicts}.

\subsection*{Retrieval}

\begin{description}[leftmargin=3.2em,style=nextline]
  \item[\textbf{H1.1} (RQ1.1)] Each formula index, added to the strongest text-only retrieval,
    improves answer retrieval over that retrieval alone.
    \hfill Section~\ref{sec:res-add-one}
  \item[\textbf{H1.2} (RQ1.2)] Adding all three formula indexes improves answer retrieval.
    \textbf{(a)}~The three-index system improves on the strongest text-only retrieval.
    \textbf{(b)}~It improves on adding one.
    \hfill Section~\ref{sec:res-add-three}
  \item[\textbf{H1.3} (RQ1.2)] How the indexes are merged makes a difference: merging the three
    formula indexes into one vote before fusing, \emph{blocked}, beats letting each of them vote
    separately, \emph{flat}.
    \hfill Section~\ref{sec:res-fusion}
  \item[\textbf{H1.4} (RQ1.2)] Every index contributes: removing any one of them makes the fused
    system worse.
    \hfill Section~\ref{sec:res-ablation}
\end{description}

\subsection*{Generation}

\begin{description}[leftmargin=3.2em,style=nextline]
  \item[\textbf{H2.1} (RQ2.1)] With the prompt held fixed, the multi-index retriever's context
    makes answers more factually correct than generation without retrieval.
    \hfill Section~\ref{sec:gen-res-context}
  \item[\textbf{H2.2} (RQ2.1)] The multi-index retriever's context makes answers both more
    factually correct and more faithful than the strongest text-only retrieval's context.
    \hfill Section~\ref{sec:gen-res-context}
  \item[\textbf{H2.3} (RQ2.2)] The better the retrieved context, the more the answer relies on it:
    faithfulness rises with the quality of the retrieved context across configurations.
    \hfill Section~\ref{sec:gen-res-tau}
  \item[\textbf{H2.4} (RQ2.2)] The better the retrieved context, the more correct the answer:
    factual correctness rises with the quality of the retrieved context across configurations.
    \hfill Section~\ref{sec:gen-res-tau}
  \item[\textbf{H3.1} (RQ3)] The conclusions of RQ2 are not specific to the smallest generator.
    \textbf{(a)}~Each comparison RQ2.1 makes keeps its direction at 3~billion and 7~billion
    parameters. \textbf{(b)}~Factual correctness and faithfulness still rise with retrieval quality
    at both sizes.
    \hfill Section~\ref{sec:gen-res-sizes}
\end{description}

% --- §1.5 and §1.6 written 2026-09-27.  Drafted in the chat as prose and approved before this
%   write (CLAUDE.md §2, steps 5-6).  Snapshot: chapter_01_introduction.pre1516.bak.tex.
%
% §1.5 -- FOUR points, Abram's selection and his own wording for two of them, polished.  The list
%   shape follows the one he gave from a fellow student's thesis: a one-line lead-in, then a
%   numbered list whose items open on a noun phrase.  ~210 words, cut from a five-point ~400-word
%   draft.  enumerate settings match §7.2 and §7.3 so the thesis's numbered lists read alike.
%
% WHAT THE FOUR POINTS CLAIM, against manifest §1d:
%   1 = C1, the architecture and the comparison of merge arrangements.
%   2 = C3 NARROWED to the judge validation alone.  §1d's C3 also covers hole-rate reporting, the
%       prime measures and the demonstration that the scoring convention decides the apparent
%       quality of a high-hole-rate ranking (claim C-4).  Chapter 1 claims none of that half; it is
%       argued in §5.3, which stays C-4's only site.
%   3 = C4, the containment diagnostic.
%   4 = C5, REWORDED.  ⚠ §1d suspends C5 and §15 O-1 lists it as blocked; both rest on O-F, the
%       prompt asymmetry, which the v3 re-run closed (§15 O-23, 2026-09-26, and §9b's note).  Every
%       Chapter 6 figure now comes from R9, so the suspension is spent.  §1d's C5 wording -- that
%       retrieval gains attenuate below detectability at the generation stage -- is NOT what R9
%       shows and is superseded by C-7 as reworded 2026-09-27; this point therefore claims the
%       MEASUREMENT and states no outcome.  Manifest §1d and O-1 updated with this write.
%   ⚠ C2 IS CLAIMED NOWHERE IN CHAPTER 1 (Abram, 2026-09-27): the from-scratch formula stack and its
%       check against a published formula-search system.  It survives as §2.2's closing sentence and
%       as §5.3's component guard, so nothing dangles, but the contributions list no longer carries
%       one of the two Exposé contributions.
%
% BANS HONOURED, each of which changed a drafted sentence:
%   - Point 1 does NOT ask whether the merge or its components matters more.  Abram's draft wording
%     was "whether the fusion or the components of the fusion matters the most"; that is the
%     comparison §1c-i rules out (CLAUDE.md §4 names this exact sentence) and "matters most" is
%     review E4's ranking clause, swept to zero in Chapter 5.  It reads instead that both are varied
%     one at a time and both decide the outcome, which is §5.6's own conclusion.
%   - NO FIGURE anywhere in §1.5, so nothing here can go stale against the registry.
%   - NO COINED TERM: not blocked, flat, formula block, fanout, floor, configuration, judge pool,
%     prime, path tuple, SLT or OPT, and not reciprocal rank fusion -- all are defined between
%     Chapter 3 and §6.2, thirty to fifty pages later, and §15 O-20 (b) is still open because §1.1
%     and §1.2 are unwritten.  Each is described instead.  Tighten to the terms once §1.1-§1.2
%     establish them.
%   - Point 4's third sentence says the measurement is REPEATED at two larger generators as a check
%     that its answers are not specific to the smallest -- RQ3 as §1b, §6.1 and §6.3.4 state it.  The
%     drafted "so that what a context adds is separable from what the generator already holds" was
%     cut on Abram's reading: it describes a decomposition the thesis never performs.
%
% ⚠ THE EMERGENT FLAG IS NOT STATED.  Manifest §5 row 1.5 asks that the contributions not planned in
%   the proposal be flagged; three of these four are.  Offered twice in the chat as one sentence --
%   "The first of these was set out at the beginning of the work; the other three emerged from it" --
%   and not taken.  Left out; manifest §5 row 1.5 records it as unmet.  A form naming the Exposé
%   directly is ruled out anyway by the §7.2 precedent of 2026-09-27, where "and the proposal this
%   work began from asked about it" was cut for pointing the reader at a document they have not seen.
%
% ⚠ THE THREE SCOPE DECISIONS STILL HAVE NO HOME.  §1c-i and §10 expect no significance testing, no
%   comparison with published answer-retrieval systems, and retrieval reported under the judge's
%   labels alone to be stated in §1.6.  A 180-word road map cannot carry them and Chapter 1 has no
%   scope section, so with §7.2 also written without them the thesis states none of the three.
%   Raised twice in the chat and open.
\section{Contributions}\label{sec:intro-contributions}

The main contributions of this thesis are as follows.

\begin{enumerate}[leftmargin=1.7em,itemsep=0.7em,topsep=0.5em]

\item The construction of three structure-aware formula indexes beside a sparse and a dense text
index, and a systematic comparison of the ways their rankings can be combined. Two things are
varied one at a time: how the three formula rankings enter a combination, and which indexes are in
it. Both decide what the combined system retrieves.

\item The validation of a language model as a relevance judge on the collection's four-level
relevance scale. Whether its labels order systems as the assessors' labels do is tested separately
from whether it agrees answer by answer, and the smallest difference between two scores it can
resolve is measured and then applied wherever two systems are compared.

\item A diagnostic for whether a retrieval signal matches what a task counts as relevant. It
contrasts how far the signal separates relevant from non-relevant answers over a topic's whole
judged set with how far it separates them among the best-matching answers, where a ranking is
decided.

\item An end-to-end measurement of what a difference in retrieval reaches in the generated answer.
Every retrieval system compared here supplies a context to one generator under one prompt, beside
answers written with no context at all, so that two things can be set against retrieval quality:
how much of an answer is drawn from its context, and whether what it says is correct. The
measurement is repeated at two larger generators, as a check that its answers are not specific to
the smallest.

\end{enumerate}

% §1.6 -- 178 words against Abram's 180-word limit, one or two sentences per remaining chapter
%   (manifest §5 row 1.6).  Seven chapters in seven sentences across six paragraphs, Chapters 7 and 8
%   sharing the last.
% EVERY chapter is named by \ref against its semantic key, never as a literal number: the repo rule
%   is that no chapter or table number is written out in the source (CLAUDE.md §6), and a road map is
%   the one place a stale number would be invisible.  Keys used: ch:related-work, ch:background,
%   ch:setup, ch:retrieval, ch:generation, ch:discussion (Chapter 7, "Concluding remarks" -- the key
%   predates the rename and a key is not prose) and ch:ai-usage.
% Navigational by design, which is what a structure section is; §3's argumentative test governs
%   forward references inside a chapter, not this one.  Each sentence says what the chapter DOES, so
%   no sentence is merely a title restated: Chapter 4's names the two preconditions, Chapter 5's the
%   five moves of its argument, Chapter 6's the two measures without naming them.
% No outcome, no verdict, no hypothesis ID and no figure -- the reader meets those in their chapters.
% Deliberately NOT here: the RQ-to-chapter mapping (§1.3 carries it with its \hfill pointers) and the
%   three scope decisions (see the §1.5 note above).
\section{Structure of the thesis}\label{sec:intro-structure}

Chapter~\ref{ch:related-work} sets this work against research on mathematical information
retrieval, formula retrieval systems, retrieval-augmented generation and evaluation with incomplete
judgments, and closes on the gap these questions address.

Chapter~\ref{ch:background} defines what the rest of the thesis uses: the collection and its two
tasks, formula representations, indexing, retrieval and fusion, and the measures for a ranking and
for a generated answer.

Chapter~\ref{ch:setup} fixes what the two experimental chapters share: the topics and judgments,
the corpus as the system reads it, the system end to end, and the two preconditions.

Chapter~\ref{ch:retrieval} answers the first question: it builds the systems compared, answers two
threats to their scores, reports the results, locates where the formula signal lies, and rules on
each retrieval hypothesis.

Chapter~\ref{ch:generation} answers the second and the third, generating answers from those
systems' contexts and scoring each for its correctness and for how much of it came from that
context.

Chapter~\ref{ch:discussion} summarises what the two answers say together, what they leave
unsettled, and where the work goes next. Chapter~\ref{ch:ai-usage} declares the use of AI tools.
